\documentclass[conference]{IEEEtran}
\IEEEoverridecommandlockouts

\usepackage{cite}
\usepackage{amsmath,amssymb,bm}
\usepackage{graphicx}
\usepackage{xcolor}
\usepackage{url}
\usepackage{booktabs}
\usepackage{multirow}
% Local: figures in runs/workshop4/figures/  Overleaf: figures/ (see scripts/pack_report.py)
\graphicspath{{../runs/workshop4/figures/}{figures/}}

% ---------------------------------------------------------------------------
% Every number quoted in the text comes from generated/results.tex, which is
% written by  python src/w4/make_report_w4.py  from the files in runs/workshop4/.
% Re-running the experiments therefore updates the report automatically.
% \figw{width}{file}{caption}{label}: full-width figure, only if the PNG exists.
% ---------------------------------------------------------------------------
% Auto-generated by src/make_report_w4.py - do not edit by hand
\newcommand{\zstar}{2}
\newcommand{\betastar}{4}
\newcommand{\epochs}{50}
\newcommand{\daenoise}{0.2}
\newcommand{\aeMSE}{0.2312}
\newcommand{\aeRsq}{0.777}
\newcommand{\aeCorr}{0.878}
\newcommand{\vaeMSE}{0.5144}
\newcommand{\vaeRsq}{0.505}
\newcommand{\vaeCorr}{0.705}
\newcommand{\aeMSEsafe}{0.2240}
\newcommand{\aeMSEfail}{0.3571}
\newcommand{\vaeMSEsafe}{0.5151}
\newcommand{\vaeMSEfail}{0.5028}
\newcommand{\aeTopAnom}{25.9\%}
\newcommand{\vaeTopAnom}{29.6\%}
\newcommand{\aeTopXdot}{0.83}
\newcommand{\vaeTopXdot}{1.06}
\newcommand{\allXdot}{0.47}
\newcommand{\aeTopX}{1.35}
\newcommand{\vaeTopX}{1.35}
\newcommand{\allX}{0.43}
\newcommand{\aeTopFar}{37.0\%}
\newcommand{\vaeTopFar}{48.1\%}
\newcommand{\allFar}{8.8\%}
\newcommand{\aeTopThetadot}{0.57}
\newcommand{\vaeTopThetadot}{0.75}
\newcommand{\allThetadot}{0.35}
\newcommand{\testAnomRate}{5.4\%}
\newcommand{\vaeKL}{1.12}
\newcommand{\vaeActive}{2}
\newcommand{\aeAUC}{0.838}
\newcommand{\aePR}{0.383}
\newcommand{\aeFone}{0.383}
\newcommand{\aeRec}{0.417}
\newcommand{\aePrec}{0.354}
\newcommand{\vaeAUC}{0.934}
\newcommand{\vaePR}{0.607}
\newcommand{\vaeFone}{0.594}
\newcommand{\vaeklAUC}{0.961}
\newcommand{\vaeklPR}{0.674}
\newcommand{\vaeklFone}{0.628}
\newcommand{\daeAUC}{0.808}
\newcommand{\daePR}{0.336}
\newcommand{\supAUC}{0.997}
\newcommand{\supPR}{0.956}
\newcommand{\supFone}{0.865}
\newcommand{\gbAUC}{0.996}
\newcommand{\gbPR}{0.935}
\newcommand{\gbFone}{0.860}
\newcommand{\lrAUC}{0.981}
\newcommand{\lrPR}{0.791}
\newcommand{\aeQrec}{0.300}
\newcommand{\aeQprec}{0.544}
\newcommand{\vaeklQrec}{0.545}
\newcommand{\vaeklQprec}{0.725}
\newcommand{\aeAllAUC}{0.686}
\newcommand{\silRaw}{0.208}
\newcommand{\silAE}{0.219}
\newcommand{\silVAE}{0.210}
\newcommand{\silSup}{0.372}
\newcommand{\probeRaw}{0.981}
\newcommand{\probeAE}{0.924}
\newcommand{\probeVAE}{0.571}
\newcommand{\probeSup}{0.997}
\newcommand{\probePRraw}{0.791}
\newcommand{\probePRAE}{0.377}
\newcommand{\probePRVAE}{0.065}
\newcommand{\probePRSup}{0.952}
\newcommand{\silMlp}{0.647}
\newcommand{\probeMlp}{0.997}
\newcommand{\probePRMlp}{0.961}
\newcommand{\dimSup}{64}
\newcommand{\aeNoisyMSE}{0.3086}
\newcommand{\daeNoisyMSE}{0.2873}
\newcommand{\aeCleanMSE}{0.2766}
\newcommand{\daeCleanMSE}{0.2798}
\newcommand{\noiseVar}{0.0400}
\newcommand{\genWvae}{0.287}
\newcommand{\genWae}{0.291}
\newcommand{\genConsVae}{0.024}
\newcommand{\genConsAe}{0.011}
\newcommand{\genBeyond}{0.1\%}
\newcommand{\activeHalf}{2}
\newcommand{\klHalf}{4.37}
\newcommand{\mseHalf}{0.2394}
\newcommand{\silbetaHalf}{0.148}
\newcommand{\activeOne}{2}
\newcommand{\klOne}{3.29}
\newcommand{\mseOne}{0.2707}
\newcommand{\silbetaOne}{0.143}
\newcommand{\activeFour}{2}
\newcommand{\klFour}{1.12}
\newcommand{\mseFour}{0.5144}
\newcommand{\silbetaFour}{0.210}
\newcommand{\trainRowsAll}{24,782}
\newcommand{\trainRowsNormal}{23,402}
\newcommand{\aucbetaHalf}{0.833}
\newcommand{\aucbetaOne}{0.804}
\newcommand{\aucbetaFour}{0.934}
\newcommand{\vaeAllAUC}{0.389}
\newcommand{\errZeroNormal}{1.076}
\newcommand{\errTenNormal}{0.150}
\newcommand{\errZeroAll}{0.206}
\newcommand{\errTenAll}{0.177}
\newcommand{\aeRsqzTwo}{0.777}
\newcommand{\aeMSEzTwo}{0.2312}
\newcommand{\aeAUCzTwo}{0.838}
\newcommand{\aeRsqzFour}{0.948}
\newcommand{\aeMSEzFour}{0.0539}
\newcommand{\aeAUCzFour}{0.777}
\newcommand{\aeRsqzEight}{0.997}
\newcommand{\aeMSEzEight}{0.0036}
\newcommand{\aeAUCzEight}{0.828}
\newcommand{\monVaeAUC}{0.897}
\newcommand{\monVaePR}{0.182}
\newcommand{\monVaeCaught}{94\%}
\newcommand{\monVaeFalse}{11.9\%}
\newcommand{\alarmDqnVae}{23.6\%}
\newcommand{\alarmPdVae}{1.3\%}
\newcommand{\alarmRandVae}{14.3\%}
\newcommand{\monWoneAUC}{0.987}
\newcommand{\monWonePR}{0.755}
\newcommand{\monWoneCaught}{100\%}
\newcommand{\monWoneFalse}{4.6\%}
\newcommand{\alarmDqnWone}{9.6\%}
\newcommand{\alarmPdWone}{0.0\%}
\newcommand{\alarmRandWone}{23.7\%}
\newcommand{\monWtwoAUC}{0.981}
\newcommand{\monWtwoPR}{0.691}
\newcommand{\monWtwoCaught}{92\%}
\newcommand{\monWtwoFalse}{4.5\%}
\newcommand{\alarmDqnWtwo}{9.5\%}
\newcommand{\alarmPdWtwo}{0.0\%}
\newcommand{\alarmRandWtwo}{24.6\%}
\newcommand{\monStates}{202,764}
\newcommand{\monPreRate}{1.7\%}
\newcommand{\monFailedEps}{690}
\newcommand{\dqnFailNominal}{3}
\newcommand{\dqnFailNoisy}{51}
\newcommand{\dqnTheta}{6.1}
\newcommand{\pdTheta}{0.2}
\newcommand{\dataTheta}{3.3}

% Allow up to two full-width figures at the top of a page (many figures in this report)
\renewcommand{\dbltopfraction}{0.9}
\renewcommand{\dblfloatpagefraction}{0.8}
\renewcommand{\topfraction}{0.9}
\renewcommand{\textfraction}{0.1}
\setcounter{dbltopnumber}{2}
\newcommand{\figw}[4]{%
  \begin{figure*}[!t]
    \centering
    \includegraphics[width=#1]{#2.png}
    \caption{#3}
    \label{#4}
  \end{figure*}}

\begin{document}

\title{Autoencoders and Variational Autoencoders for Unsupervised Failure Detection in CartPole Trajectories}

\author{\IEEEauthorblockN{Miguel Angel Torres Ramirez}
\IEEEauthorblockA{Dept. of Computer Engineering\\
Universidad Distrital Francisco Jos\'e de Caldas\\
Email: migatorresr@udistrital.edu.co}
\and
\IEEEauthorblockN{Cristian David Romero Gil}
\IEEEauthorblockA{Dept. of Computer Engineering\\
Universidad Distrital Francisco Jos\'e de Caldas\\
Email: cdromerog@udistrital.edu.co}
\and
\IEEEauthorblockN{Omar Esteban Pe\~na Sanchez}
\IEEEauthorblockA{Dept. of Computer Engineering\\
Universidad Distrital Francisco Jos\'e de Caldas\\
Email: oepenas@udistrital.edu.co}}

\maketitle

\begin{abstract}
This report studies what unsupervised representation learning reveals about the CartPole
failure-prediction dataset introduced in Workshop~1. We train fully-connected
autoencoders (AE) with latent sizes $z\in\{2,4,8\}$ and variational autoencoders (VAE)
with $\beta\in\{0.5,1,4\}$ on the same 16 engineered features and the same episode-level
split, without using the failure labels. Pre-failure states are treated as anomalies:
models trained only on safe states score each test state by its reconstruction error.
Accurate reconstruction does not imply good detection. The AE with $z=8$ reconstructs
almost perfectly ($R^2=\aeRsqzEight$), yet the selected detectors use the smallest code
($z=\zstar$) and the strongest regularisation ($\beta=\betastar$). The AE reaches ROC-AUC
$=\aeAUC$. The VAE reaches \vaeAUC{} with its reconstruction error and \vaeklAUC{} when the KL
term is added (PR-AUC $=\vaeklPR$), all without labels. Supervised models remain clearly
superior (Workshop~2 LSTM: ROC-AUC $=\supAUC$, PR-AUC $=\supPR$). The latent-space analysis
explains this gap. Unsupervised codes are organised by the dominant physical variables, and
failure states lie at the periphery of the latent space rather than along a separable
direction. The frozen Workshop~2 LSTM, in contrast, organises its hidden state around
failure risk (balanced silhouette \silSup{} versus \silAE{} for the AE and \silVAE{} for
the VAE), and a same-architecture control shows that the labels, not the architecture, make
the difference. Used as a label-free monitor of the Workshop~3 controllers, the VAE detects
pre-failure states on unseen policies (ROC-AUC $=\monVaeAUC$) and flags the DQN agent's
unusual operating regime.
\end{abstract}

\begin{IEEEkeywords}
autoencoder, variational autoencoder, $\beta$-VAE, anomaly detection, representation learning, latent space, CartPole.
\end{IEEEkeywords}

%% ======================================================================
\section{Introduction}
Workshop~1 framed failure prediction in CartPole \cite{barto1983,brockman2016} as a
supervised problem: a classifier receives the state at time $t$ and must raise an alarm
during the last five steps before the pole falls or the cart leaves the track. Supervision,
however, has a cost. Failures are rare (5.5\% of the states), they must be
labelled after the fact, and a monitoring system deployed on a new controller may never have
seen one. \emph{Unsupervised} learning offers an alternative: learn what normal operation
looks like and flag whatever does not fit.

Autoencoders learn a compressed code by reconstructing their input
\cite{hinton2006,goodfellow2016}; variational autoencoders add a probabilistic latent space
regularised towards a prior \cite{kingma2014,doersch2016}. Both give two tools for this
domain: a latent space to inspect, and a reconstruction error that can serve as an anomaly
score \cite{chalapathy2019}. The central question of this workshop is: \emph{what structure
do unsupervised models reveal about the dataset, and how can these representations enhance
our understanding and the downstream task?}

The contributions are:
(i) an AE and a $\beta$-VAE adapted to tabular control data, with a latent-size sweep and a
$\beta$ sweep;
(ii) a reconstruction-based failure detector trained only on safe states, compared with the
Workshop~1 baselines and the Workshop~2 deep LSTM;
(iii) a latent-space analysis (t-SNE, UMAP, silhouette, linear probes and per-dimension
correlations) contrasting the unsupervised codes with the frozen Workshop~2 encoder; and
(iv) three extensions: a denoising AE, an evaluation of the VAE as a generative model, and
the use of the VAE as a label-free monitor of the Workshop~3 reinforcement-learning
controllers.

%% ======================================================================
\section{Data and Problem Setup}

\subsection{Dataset and Inputs}
We reuse the Workshop~1 dataset without changes: 35,874 states from 400 CartPole-v1
episodes, generated with a mixed noisy-expert/random policy. Each input is the vector of the
16 engineered features of Workshop~1. These are the raw state $(x,\dot{x},\theta,\dot\theta)$
and the action, the distances to the failure limits, first differences, 3-step rolling means
and time-to-limit estimates. None of them uses future information. Features are z-scored
with statistics from the training episodes only. The episode-level 70/15/15 split stored in
\texttt{data/splits.json} is reused, so every result is directly comparable with
Workshop~1 and there is no leakage between consecutive states.

\subsection{Normal and Anomalous States}
The dataset has a natural normal/abnormal partition: a state is \emph{normal} if
\texttt{label}$=0$ (safe) and \emph{anomalous} if \texttt{label}$=1$ (failure within five
steps). The test set contains 5,330 states, \testAnomRate{} of them anomalous. Each
unsupervised model is trained in two regimes:
\begin{itemize}
\item \textbf{All rows} (\trainRowsAll{} training states). These models are used to measure
reconstruction quality and to analyse the latent space.
\item \textbf{Normal rows only} (\trainRowsNormal{} training states). These models are used
for anomaly detection, as the workshop prescribes.
\end{itemize}
Labels are used only to evaluate the models, to tune thresholds on validation and to select
hyper-parameters on validation. They are never used to train an unsupervised model.

%% ======================================================================
\section{Methods}

\subsection{Autoencoder}
Because the data are tabular, encoder and decoder are fully-connected networks
(Table~\ref{tab:arch}). The decoder mirrors the encoder and ends in a \emph{linear} layer,
since the targets are unbounded z-scores and a Tanh or Sigmoid would clip them. The loss is
the mean squared error (MSE). The workshop suggests $z\in\{32,64,128\}$ and a latent size of
about 5--10\% of the input dimension. With 16 inputs, the first suggestion would give an
\emph{over-complete} code that can learn the identity, while 5--10\% means 1--2 units. We
therefore sweep $z\in\{2,4,8\}$, which brackets the guideline and stays below the input
dimension.

\begin{table}[t]
\caption{Network Architectures (ReLU Hidden Activations)}
\label{tab:arch}
\centering
\resizebox{\columnwidth}{!}{%
\begin{tabular}{ll}
\hline
\textbf{Block} & \textbf{Layers} \\ \hline
AE encoder & $16 \rightarrow 64 \rightarrow 32 \rightarrow z$ \\
VAE encoder & $16 \rightarrow 64 \rightarrow 32 \rightarrow (\bm\mu, \log\bm\sigma^2)\in\mathbb{R}^{z}\times\mathbb{R}^{z}$ \\
Decoder (both) & $z \rightarrow 32 \rightarrow 64 \rightarrow 16$ (linear output) \\
W2 LSTM (frozen) & window $16\times5$ (raw) $\rightarrow$ LSTM ($2\times64$) $\rightarrow 32 \rightarrow 1$ \\
MLP control & $16 \rightarrow 64 \rightarrow 32 \rightarrow z$ (ReLU) $\rightarrow 1$ (logit) \\
\hline
\end{tabular}}
\end{table}

\subsection{Variational Autoencoder}
The VAE encoder outputs the parameters of a diagonal Gaussian posterior
$q_\phi(\mathbf z|\mathbf x)=\mathcal N(\bm\mu,\mathrm{diag}\,\bm\sigma^2)$. A latent sample
is obtained with the \emph{reparameterisation trick}
\begin{equation}
\mathbf z=\bm\mu+\bm\sigma\odot\bm\epsilon,\qquad \bm\epsilon\sim\mathcal N(\mathbf 0,\mathbf I),
\end{equation}
which moves the randomness to an external input. The sample then becomes a deterministic,
differentiable function of $\bm\mu$ and $\bm\sigma$, so gradients flow through the encoder.
The network minimises the negative evidence lower bound (ELBO) with a weight $\beta$ on the
KL term \cite{higgins2017}:
\begin{align}
\mathcal L &= \underbrace{\textstyle\sum_{j=1}^{16}(x_j-\hat x_j)^2}_{\mathcal L_{\text{recon}}}
 + \beta\,\mathcal L_{\text{KL}}, \\
\mathcal L_{\text{KL}} &= -\tfrac12\textstyle\sum_{k=1}^{z}\left(1+\log\sigma_k^2-\mu_k^2-\sigma_k^2\right).
\end{align}
The reconstruction term is the squared error \emph{summed} over features, which corresponds
to a Gaussian decoder with fixed variance. This keeps it on the same per-sample scale as the
KL term, so $\beta$ has a clear meaning. To prevent posterior collapse (KL vanishing), we use
KL annealing \cite{bowman2016}: $\beta$ rises linearly from 0 to its target over the first
10 epochs. At evaluation time the VAE encodes and reconstructs with the posterior mean
$\bm\mu$. Collapse is diagnosed with the number of \emph{active units}: dimensions whose
$\mathrm{Var}_{\mathbf x}[\mu_k]$ exceeds 0.01 \cite{burda2016}.

\subsection{Denoising Autoencoder (Extension)}
Following \cite{vincent2010}, a denoising AE with the selected latent size receives
$\tilde{\mathbf x}=\mathbf x+\mathcal N(0,\daenoise^2\mathbf I)$ and is trained to reconstruct
the clean $\mathbf x$. The noise acts as a regulariser against the identity mapping.

\subsection{Workshop~2 Model and Supervised Control}
The supervised reference is the Workshop~2 network itself, loaded from its checkpoint and
\emph{not} retrained. It is a two-layer LSTM with 64 hidden units that reads a past-only
window of the last 16 raw steps $(x,\dot x,\theta,\dot\theta,a)$, followed by an MLP head
(Table~\ref{tab:arch}). It uses the same episode split, so its outputs align row by row with
the autoencoder inputs. Re-running it reproduces the Workshop~2 test results exactly. Its
final hidden state ($\dimSup$-D) is the \emph{frozen supervised encoder} used in the
latent-space comparison, and its probabilities give the Workshop~2 row of the cumulative
table.

The LSTM differs from the autoencoders in two ways: it is trained with labels, and it sees
16 steps of history instead of one engineered vector. To separate these two effects, we also
train a \emph{supervised control}: an MLP with the same input, the same trunk and the same
bottleneck width $z$ as the AE, trained with binary cross-entropy weighted by the class ratio.
It appears only in the latent-space table, marked with $\ddagger$.

\subsection{Training and Model Selection}
All models use Adam \cite{kingma2015} (learning rate $10^{-3}$, batch 64) for \epochs{}
epochs with seed 42. We log train and validation losses every epoch and keep the best
validation epoch. For the VAE, the best epoch is chosen with the target-$\beta$ objective and
only after the warm-up. The supervised control keeps the epoch with the best validation
PR-AUC.

Selection uses validation data only. First, $z^\star$ is the latent size whose
normal-rows AE has the highest validation anomaly ROC-AUC. Second, the VAE uses $z^\star$,
and $\beta^\star$ is selected the same way. This gives $z^\star=\zstar$ and
$\beta^\star=\betastar$. The \emph{best} checkpoints (\texttt{ae\_best.pt},
\texttt{vae\_best.pt}) are the all-rows models with these settings.

\subsection{Evaluation Protocol}
\textbf{Reconstruction.} On the test set we report MSE in standardised units, the explained
variance $R^2$ and the mean Pearson correlation $\bar\rho$ between each feature and its
reconstruction. These are the continuous-data metrics the workshop suggests; SSIM and PSNR
do not apply to tabular data.

\textbf{Anomaly detection.} The score is the per-state reconstruction MSE. For the VAE we
also use $\mathcal L_{\text{recon}}+\mathcal L_{\text{KL}}$. We report the threshold-free
ROC-AUC and PR-AUC. Precision, recall and F1 are reported at the threshold that maximises F1
on \emph{validation}, the same rule as in Workshop~1. As a label-free alternative, we also
use the 99th percentile of the validation normal scores.

\textbf{Latent space.} Test states are encoded (with $\bm\mu$ for the VAE) and compared with
the input features, the frozen Workshop~2 encoder and the supervised control. We use four
measurements:
(a) the silhouette score \cite{rousseeuw1987} of the label partition, after z-scoring each
dimension. It is computed on all rows and, because of the 17:1 imbalance, on ten
class-balanced subsamples.
(b) A \emph{linear probe} \cite{alain2016}: a logistic regression fitted on training codes
and evaluated on test, which measures how much failure information is linearly available.
(c) 2-D t-SNE \cite{vandermaaten2008} and UMAP \cite{mcinnes2018} maps.
(d) The correlation between each latent dimension and the physical variables.

%% ======================================================================
\section{Results}

\figw{\textwidth}{fig-w4-training-curves}{Training curves. (a) AE reconstruction MSE for each
latent size. (b) VAE reconstruction MSE and (c) validation KL divergence for each $\beta$
($z=z^\star$); the shaded band is the KL warm-up.}{fig:curves}

\subsection{Training Dynamics}
Fig.~\ref{fig:curves} shows stable training without overfitting: the train and validation
curves stay close in every run. The AE error falls sharply with the latent size (test MSE
\aeMSEzTwo, \aeMSEzFour{} and \aeMSEzEight{} for $z=2,4,8$). With $z=8$ the code is only
half the input width, yet the reconstruction is nearly perfect ($R^2=\aeRsqzEight$), because
the 16 engineered features are highly redundant. Several of them are deterministic functions
of the raw state, for example $|\theta|$ and $0.2095-|\theta|$. Even so, the validation
anomaly criterion selects the \emph{smallest} code, $z^\star=\zstar$
(Table~\ref{tab:sweep}). This first sign that reconstruction accuracy and detection ability
diverge is examined in Section~\ref{sec:anomaly}.

The VAE curves show the annealing mechanism at work. For $\beta=0.5$ and $\beta=1$ the
reconstruction error keeps decreasing throughout training. For $\beta=4$ the error rises
during the warm-up, as the KL weight grows, and settles at a higher level once the weight
reaches its target. The KL term then stays at \klFour{} nats per sample, compared with
\klHalf{} nats for $\beta=0.5$.

\subsection{Reconstruction Quality}
Table~\ref{tab:recon} summarises the test reconstruction. The selected AE ($z=\zstar$)
reaches MSE $=\aeMSE$ and $\bar\rho=\aeCorr$. The selected VAE ($\beta=\betastar$) is less
accurate (MSE $=\vaeMSE$, $R^2=\vaeRsq$). This is the expected rate-distortion trade-off: as
$\beta$ grows, every nat of information in $\mathbf z$ costs more, so the model keeps only the
dominant factors. Both latent dimensions remain active for every $\beta$
(Table~\ref{tab:recon}), so there is no posterior collapse. With a two-unit code, $\beta$
reduces how much information each dimension carries rather than how many dimensions survive.

\begin{table}[t]
\caption{Test Reconstruction Quality (All-Rows Models). MSE$_0$/MSE$_1$: Safe/Failure States; AU: Active Units; $\star$: Selected}
\label{tab:recon}
\centering
\setlength{\tabcolsep}{3.2pt}
\resizebox{\columnwidth}{!}{\begin{tabular}{lccccccc}
\hline
\textbf{Model} & \textbf{MSE} & $\mathbf{R^2}$ & $\bar{\boldsymbol{\rho}}$ & \textbf{MSE$_0$} & \textbf{MSE$_1$} & \textbf{KL} & \textbf{AU} \\ \hline
AE, $z=2$$^\star$ & 0.2312 & 0.777 & 0.878 & 0.2240 & 0.3571 & -- & -- \\
AE, $z=4$ & 0.0539 & 0.948 & 0.974 & 0.0529 & 0.0716 & -- & -- \\
AE, $z=8$ & \textbf{0.0036} & \textbf{0.997} & \textbf{0.998} & 0.0036 & 0.0033 & -- & -- \\
\hline
VAE, $z=2$, $\beta=0.5$ & 0.2394 & 0.770 & 0.870 & 0.2318 & 0.3719 & 4.37 & 2/2 \\
VAE, $z=2$, $\beta=1$ & 0.2707 & 0.739 & 0.852 & 0.2645 & 0.3777 & 3.29 & 2/2 \\
VAE, $z=2$, $\beta=4$$^\star$ & 0.5144 & 0.505 & 0.705 & 0.5151 & 0.5028 & 1.12 & 2/2 \\
\hline
\end{tabular}
}
\end{table}

\figw{\textwidth}{fig-w4-reconstructions}{Sample reconstructions of eight test states (four safe,
four within five steps of failure). Columns are the 16 standardised features. From left to
right: original, AE reconstruction ($z=z^\star$), VAE reconstruction ($\beta=\beta^\star$)
and element-wise absolute-error maps on a common scale.}{fig:samples}

Fig.~\ref{fig:samples} shows what the two models keep. Both reproduce the dominant pattern
of each state: the sign and size of $\theta$ and $\dot\theta$ and the distance-to-limit
features. The pre-failure rows have a strong and coherent pattern of this kind, so they are
reconstructed well. The largest errors are on the binary action, the cart velocity and the
change in angular velocity (Fig.~\ref{fig:errors}a). The VAE also pulls the safe states
towards the mean, which shows as pale reconstructions in the upper rows.

\figw{\textwidth}{fig-w4-error-analysis}{Error analysis on the test set. (a) MSE per feature.
(b) Median reconstruction error versus the number of steps left before a failure, for the
AE trained on all rows and on normal rows only. (c) Distribution of the per-state error by
class.}{fig:errors}

\textbf{Which states are hardest to reconstruct?} The hardest states are the rare, extreme
configurations of the system. Among the 1\% of test states with the largest AE error,
\aeTopFar{} have the cart more than 1.5~m from the centre (\allFar{} overall), and their mean
$|\dot x|$ is \aeTopXdot{}~m/s (\allXdot{}~m/s overall). Pre-failure states are
over-represented among them (\aeTopAnom{} versus a \testAnomRate{} base rate), but most of
the hardest states are safe. The same holds for the VAE (\vaeTopFar{} far from the centre,
\vaeTopAnom{} pre-failure). Fig.~\ref{fig:errors}b shows how training data shapes this
pattern. For the AE trained on safe states only, the median error grows from \errTenNormal{}
ten steps before a failure to \errZeroNormal{} at the failing step. For the AE trained on all
rows, it barely changes (\errTenAll{} to \errZeroAll), because that model has learned to
reconstruct pre-failure states.

\subsection{Anomaly Detection}
\label{sec:anomaly}
Table~\ref{tab:anomaly} is the cumulative comparison requested by the workshop: the
Workshop~1 baselines, the Workshop~2 LSTM and the Workshop~4 detectors on the
same test episodes. Table~\ref{tab:sweep} reports the sweep over latent size, $\beta$ and
training regime. Fig.~\ref{fig:anomaly} shows the curves.

\begin{table}[t]
\caption{Cumulative Test Results for Failure / Anomaly Detection (Threshold Selected on Validation). Bold: Best Unsupervised}
\label{tab:anomaly}
\centering
\setlength{\tabcolsep}{4pt}
\begin{tabular}{llccccc}
\hline
& \textbf{Model} & \textbf{ROC} & \textbf{PR} & \textbf{P} & \textbf{R} & \textbf{F1} \\ \hline
W1 & Logistic Reg. & 0.981 & 0.791 & 0.669 & 0.703 & 0.686 \\
W1 & Random Forest & 0.994 & 0.923 & 0.812 & 0.852 & 0.832 \\
W1 & Gradient Boost. & 0.996 & 0.935 & 0.834 & 0.886 & 0.860 \\
W2 & Deep LSTM & 0.997 & 0.956 & 0.812 & 0.924 & 0.865 \\
\hline
W4 & AE & 0.838 & 0.383 & 0.354 & 0.417 & 0.383 \\
W4 & VAE: recon. & 0.934 & 0.607 & \textbf{0.664} & 0.538 & 0.594 \\
W4 & VAE: recon.\,+\,KL & \textbf{0.961} & \textbf{0.674} & 0.646 & \textbf{0.610} & \textbf{0.628} \\
W4 & Denoising AE & 0.808 & 0.336 & 0.304 & 0.317 & 0.310 \\
\hline
\end{tabular}

\end{table}

\begin{table}[t]
\caption{Anomaly ROC-AUC by Latent Size, $\beta$ and Training Regime (Test in Parentheses: VAE Recon.\,+\,KL Score)}
\label{tab:sweep}
\centering
\setlength{\tabcolsep}{5pt}
\begin{tabular}{lccc}
\hline
& \multicolumn{2}{c}{\textbf{Normal rows only}} & \textbf{All rows} \\
\textbf{Model} & \textbf{Val} & \textbf{Test} & \textbf{Test} \\ \hline
AE, $z=2$$^\star$ & 0.834 & 0.838 & 0.686 \\
AE, $z=4$ & 0.744 & 0.777 & 0.609 \\
AE, $z=8$ & 0.802 & 0.828 & 0.644 \\
\hline
VAE, $\beta=0.5$ & 0.784 & 0.833 (0.914) & 0.679 \\
VAE, $\beta=1$ & 0.760 & 0.804 (0.917) & 0.615 \\
VAE, $\beta=4$$^\star$ & 0.928 & 0.934 (0.961) & 0.389 \\
\hline
\end{tabular}

\end{table}

\figw{\textwidth}{fig-w4-anomaly}{Anomaly detection on the test set. (a) ROC and (b)
precision-recall curves of the unsupervised detectors (trained on safe states only) and
the Workshop~2 LSTM; the grey line in (b) is the prevalence. (c) Distribution of the AE
anomaly score by class.}{fig:anomaly}

Four findings stand out.
\begin{enumerate}
\item \textbf{Training on normal rows is essential.} The same AE trained on all rows reaches
only ROC-AUC $=\aeAllAUC$, compared with \aeAUC{} when trained on safe states. The selected
VAE trained on all rows falls to \vaeAllAUC, below chance (Table~\ref{tab:sweep}). Once
pre-failure states are in the training data, the models reconstruct them as well as anything
else.
\item \textbf{Regularisation beats reconstruction accuracy.} The VAE reconstructs worse than
the AE but detects better (ROC-AUC \vaeAUC{} versus \aeAUC). Among the AEs, the
near-perfect $z=8$ model is not the best detector either. A strongly constrained model can
only represent the bulk of the normal data, so unusual states are pulled towards typical ones
and leave a large residual. The denoising AE did not help in this setting (\daeAUC). A
plausible explanation is that, with a two-unit code, the bottleneck already limits capacity,
and the input noise mainly blurs the extreme states that signal a failure.
\item \textbf{The KL term is informative.} Adding the KL divergence to the VAE score raises
ROC-AUC to \vaeklAUC{} and PR-AUC to \vaeklPR. Pre-failure states are encoded far from the
prior, so the posterior itself signals novelty, consistent with \cite{schlegl2017}.
\item \textbf{Supervision still wins by a wide margin.} The best unsupervised detector ranks
states well, but its PR-AUC (\vaeklPR) is clearly below Gradient Boosting (\gbPR) and the
Workshop~2 LSTM (\supPR). Its F1 at the validation threshold is \vaeklFone{}, compared with
\gbFone{} and \supFone. The score distributions overlap (Fig.~\ref{fig:anomaly}c), because many safe
states, such as large but recoverable angles under the random policy, are as unusual as the
last steps before a failure. With the label-free 99th-percentile threshold, the VAE score
has precision \vaeklQprec{} and recall \vaeklQrec.
\end{enumerate}

\subsection{Latent Space Structure}

\figw{\textwidth}{fig-w4-latent-tsne}{t-SNE maps of the test states in the input space, the AE code,
the VAE posterior mean and the frozen Workshop~2 LSTM state. Top: coloured by label. Bottom:
coloured by $|\theta|$.}{fig:tsne}

\figw{\textwidth}{fig-w4-latent-umap}{UMAP maps of the same representations as Fig.~\ref{fig:tsne}.}{fig:umap}

\begin{table}[t]
\caption{Separability of the Failure Class in Each Representation (Test). Silhouette on the Full Set, on Balanced Subsamples and on the 2-D Maps; Linear Probe ROC-AUC and PR-AUC. $\ddagger$: Supervised Control (Same Input and Trunk as the AE)}
\label{tab:latent}
\centering
\setlength{\tabcolsep}{3pt}
\resizebox{\columnwidth}{!}{\begin{tabular}{lccccccc}
\hline
& & \multicolumn{4}{c}{\textbf{Silhouette}} & \multicolumn{2}{c}{\textbf{Linear probe}} \\
\textbf{Space} & \textbf{Dim} & \textbf{All} & \textbf{Bal.} & \textbf{t-SNE} & \textbf{UMAP} & \textbf{ROC} & \textbf{PR} \\ \hline
Input features & 16 & 0.336 & 0.208 & -0.117 & -0.158 & 0.981 & 0.791 \\
AE code & 2 & 0.467 & 0.219 & 0.225 & 0.222 & 0.924 & 0.377 \\
VAE $\mu$ & 2 & 0.509 & 0.210 & 0.153 & 0.184 & 0.571 & 0.065 \\
W2 LSTM state & 64 & 0.613 & 0.372 & 0.191 & 0.127 & 0.997 & 0.952 \\
MLP control$^\ddagger$ & 2 & 0.664 & 0.647 & -- & -- & 0.997 & 0.961 \\
\hline
\end{tabular}
}
\end{table}

\figw{0.85\textwidth}{fig-w4-latent-correlation}{Absolute Pearson correlation between each latent
dimension and interpretable variables (test set). The 64-D LSTM state is summarised by its
first two principal components.}{fig:corr}

\figw{0.85\textwidth}{fig-w4-beta}{Effect of $\beta$ on the VAE ($z=z^\star$, all rows).
(a) Reconstruction MSE versus KL (rate-distortion). (b) KL per latent dimension.
(c) Balanced silhouette of the failure class and anomaly ROC-AUC of the normal-rows model.}{fig:beta}

\textbf{What do the autoencoders discover without labels?} In the input space, the t-SNE
map (Fig.~\ref{fig:tsne}) is fragmented into many small clusters of similar trajectory
segments. The failure states form several separate islands. The AE code arranges the states
along continuous filaments ordered by $|\theta|$, and the failure states sit at the
\emph{ends} of two of them, the most extreme configurations. The VAE map shows the same
pattern: failure states are found at the edges of the cloud, in several groups. UMAP gives
the same picture (Fig.~\ref{fig:umap}).

The correlation analysis (Fig.~\ref{fig:corr}) makes this concrete. Each AE and VAE dimension
mixes the cart position, the angle and their velocities, and correlates only weakly with the
label. The autoencoders find the \emph{physical} structure of the system, organised by the
variables that carry the most variance. They do not single out the direction that decides
failure. As a result, the balanced silhouette of the failure class is about the same in the
AE code (\silAE), the VAE code (\silVAE) and the input space (\silRaw).

The $\beta$ study (Fig.~\ref{fig:beta}) shows that $\beta$ acts through the amount of
information in the code. The KL drops from \klHalf{} to \klOne{} and \klFour{} nats as
$\beta$ increases from 0.5 to 1 and 4, while both dimensions stay active. The strongest
regularisation ($\beta=4$) gives both the highest silhouette (\silbetaFour, compared with
\silbetaHalf{} and \silbetaOne) and the best anomaly AUC (\aucbetaFour, compared with
\aucbetaHalf{} and \aucbetaOne). The relation is not monotonic between $\beta=0.5$ and
$\beta=1$, however, so with a single seed it should be read as a trend.

\subsection{Supervised versus Unsupervised Representations}
The frozen Workshop~2 encoder is organised differently. In its t-SNE map
(Fig.~\ref{fig:tsne}) the failure states form two compact clusters at the edge of the cloud,
next to the high-$|\theta|$ regions. Its balanced silhouette (\silSup) is clearly higher
than that of the AE (\silAE), the VAE (\silVAE) and the input space (\silRaw)
(Table~\ref{tab:latent}). Its two main components correlate with the label and with
$\theta$ and $|\theta|$, and almost not at all with $|x|$ or the action
(Fig.~\ref{fig:corr}). The autoencoders, in contrast, devote a whole dimension to the cart
position. In effect, the labels told the network which variables matter and allowed it to
discard the rest.

The supervised control isolates the role of the labels. With exactly the same input and
trunk as the AE, adding supervision raises the balanced silhouette from \silAE{} to
\silMlp. The control is even more separable than the LSTM state, because its 2-D bottleneck
keeps little besides risk, whereas the 64-D LSTM state also stores the recent dynamics of
the system. The gap between supervised and unsupervised representations is therefore caused
by the training signal, not by the architecture or the temporal context.

The linear probes quantify the value of each representation for the downstream task. A
logistic regression reaches ROC-AUC \probeSup{} on the Workshop~2 LSTM state (PR-AUC
\probePRSup), \probeMlp{} on the supervised control, \probeRaw{} on the raw features (the
Workshop~1 logistic baseline) and \probeAE{} on the AE code. On the VAE
code it reaches only \probeVAE{} (PR-AUC \probePRVAE), even though the same model is the best
unsupervised detector. The two facts are consistent with the geometry seen in
Fig.~\ref{fig:tsne}. Failure states are encoded far from the centre of the latent space but
in different directions, so a single hyperplane cannot separate them. A distance-based score,
such as the KL term, which grows with $\|\bm\mu\|^2$, or the reconstruction error, does
separate them.

Compression without labels therefore keeps failure information only in a form that suits
novelty detection, not classification. The lesson about label supervision is that labels are
not only a training signal but a \emph{relevance} signal. Unsupervised models answer ``what
varies, and what is rare?'', whereas supervised ones answer ``what matters?''. In this domain
the two questions have different answers.

\subsection{Extensions: Denoising and Generation}
\textbf{Denoising.} With inputs corrupted by Gaussian noise ($\sigma=\daenoise$), the
denoising AE reconstructs the clean test states with MSE \daeNoisyMSE{}, compared with
\aeNoisyMSE{} for the plain AE, at almost no cost on clean inputs (\daeCleanMSE{} versus
\aeCleanMSE). With a two-unit code, however, the error of the bottleneck itself is much larger
than the noise variance (\noiseVar), so neither model improves on the corrupted input as such.
Denoising becomes meaningful only with larger codes, where the identity map is a real risk.

\figw{\textwidth}{fig-w4-vae-samples}{Generative properties. Distribution of the state variables in the
test set and in 5,000 samples decoded from $\mathbf z\sim\mathcal N(\mathbf 0,\mathbf I)$ by
the VAE and, as a contrast, by the AE (physical units). Right: consistency between the
generated angle and the generated $|\theta|$ feature.}{fig:gen}

\textbf{Generation.} Sampling $\mathbf z\sim\mathcal N(\mathbf 0,\mathbf I)$ and decoding
(Fig.~\ref{fig:gen}) gives VAE states that cover the range of the real data, but they are
too concentrated around the upright equilibrium. The mean 1-Wasserstein distance over the four
state variables, divided by their standard deviation, is \genWvae{} for the VAE and \genWae{}
for the AE. The AE can be sampled this way only because its two-unit code happens to have a
similar scale, which nothing in its training enforces. The generated features follow the
physical relation $|\theta|$-feature $\approx|\theta|$ for large angles but overestimate it
near zero (mean gap \genConsVae{}~rad). Using the VAE for data augmentation would therefore
require a smaller $\beta$ or a decoder variance learned per feature.

\subsection{Extension: Label-Free Monitoring of the Workshop~3 Controllers}
Workshop~3 trained a DQN agent that balances the pole and used the Workshop~1 and~2 models
to measure how often its states look dangerous. Those detectors needed failure labels. We
repeat the measurement with the VAE trained on safe states only (reconstruction-plus-KL
score, threshold from the Workshop~4 validation episodes). This is the kind of monitor that
could be deployed on a controller whose failures have never been observed.

We follow the Workshop~3 test protocol: 100 held-out initial states, under nominal dynamics
and under actuator noise (2~N). The policies are the four Workshop~3 baselines that do not
query a failure model (random, alternating, angle rule and PD rule) and the selected DQN
agent, for \monStates{} visited states in total. A state is labelled pre-failure when its
episode ends in a failure within five steps, the dataset definition. This applies to
\monPreRate{} of the states, from \monFailedEps{} failed episodes. No model or threshold is
tuned on these trajectories. The Workshop~2 alarm rates reproduce those of the Workshop~3
report exactly. The Workshop~1 Gradient Boosting model is re-fitted by the Workshop~3 bridge,
and its rates differ slightly across library versions (for example, \alarmDqnWone{} of the
DQN states here).

\begin{table}[t]
\caption{Pre-Failure Detection on the States Visited by the Workshop~3 Policies (Pooled). Caught: Failed Episodes With an Alarm in Their Last Five Steps; False al.: Alarm Rate in Episodes That Never Fail}
\label{tab:w3}
\centering
\setlength{\tabcolsep}{4pt}
\begin{tabular}{lcccc}
\hline
\textbf{Detector} & \textbf{ROC} & \textbf{PR} & \textbf{Caught} & \textbf{False al.} \\ \hline
W4 VAE (no labels) & 0.897 & 0.182 & 94\% & 11.9\% \\
W1 Gradient Boost. & 0.987 & 0.755 & 100\% & 4.6\% \\
W2 LSTM & 0.981 & 0.691 & 92\% & 4.5\% \\
\hline
\end{tabular}

\end{table}

\figw{\textwidth}{fig-w4-w3-monitoring}{Monitoring the Workshop~3 controllers on 100 test
episodes per policy. (a, b) Share of visited states flagged by each detector; the black marks
are the true share of pre-failure states. (c) ROC curves of pre-failure detection over all
policies and both conditions.}{fig:w3}

The VAE transfers to these new state distributions (Table~\ref{tab:w3}). It reaches ROC-AUC
\monVaeAUC{} and raises an alarm in the last five steps of \monVaeCaught{} of the failed
episodes. It is clearly less precise than the supervised detectors, however. Its PR-AUC is
\monVaePR{}, compared with \monWonePR{} and \monWtwoPR{}, and it flags \monVaeFalse{} of the
states in episodes that never fail, compared with \monWoneFalse{} and \monWtwoFalse.

The alarm rates per policy (Fig.~\ref{fig:w3}) show what the VAE actually measures. The PD
rule keeps the pole almost still (mean $|\theta|=\pdTheta^\circ$) and triggers almost no
alarms (\alarmPdVae). The DQN agent fails in only \dqnFailNominal{} of 100 nominal episodes,
yet the VAE flags \alarmDqnVae{} of its states, more than the Workshop~1 (\alarmDqnWone) and
Workshop~2 (\alarmDqnWtwo) models. The agent balances the pole at a mean
$|\theta|=\dqnTheta^\circ$, almost twice the \dataTheta$^\circ$ of the dataset, a regime the
data-generating policy rarely visits. The VAE reports exactly that: this controller operates
outside the normal data distribution. Under actuator noise, the same regime turns into
failures in \dqnFailNoisy{} of 100 episodes, consistent with the fragility reported in
Workshop~3.

A label-free monitor therefore detects \emph{novelty}, not danger as such. This makes it a
useful early warning that a new controller behaves unlike anything observed before, but
labels, or the supervised models of Workshops~1 and~2, are needed to tell risky novelty
from harmless novelty.

%% ======================================================================
\section{Discussion and Limitations}
\begin{itemize}
\item \textbf{Low-dimensional tabular data.} With only 16 partly redundant inputs, even a
moderately sized code approaches the identity map. This is why the smallest code and the
strongest regularisation were selected here, unlike what is typical in image domains.
\item \textbf{Single seed and selection variance.} All results come from one seed and one
split, and some validation differences in Table~\ref{tab:sweep} are a few hundredths. The
selected $z^\star$ and $\beta^\star$ should be read as indicative rather than definitive.
\item \textbf{Different inputs.} The Workshop~2 LSTM reads 16 steps of raw history,
whereas the autoencoders read one engineered vector. The same-architecture supervised control
shows that this difference does not explain the gap, but a sequence autoencoder would make the
comparison fully symmetric.
\item \textbf{Static inputs.} Each state is encoded independently. A recurrent or windowed
autoencoder over the last few steps could capture the dynamics of an approaching failure
better.
\end{itemize}

%% ======================================================================
\section{Conclusions}
Autoencoders learn a compact description of CartPole states organised by the variables with
the most variance, not by the proximity to failure. Trained only on safe states, they turn
reconstruction error into a label-free failure detector, and the most accurate reconstructor
is not the best detector. The $\beta$-VAE with the reconstruction-plus-KL score reaches
ROC-AUC \vaeklAUC{}, whereas the AE reaches \aeAUC. Failure states lie at the periphery of the
unsupervised latent spaces, which a distance-based score exploits but a linear classifier
cannot. Supervised models remain far better (PR-AUC \supPR{} for the Workshop~2 LSTM and
\gbPR{} for Gradient Boosting, versus \vaeklPR) because labels select the variables that
matter. A same-architecture control confirms that the training signal, not the
architecture, makes the difference. In practice, unsupervised models are
useful for understanding the data and for monitoring controllers where failures have never
been observed. On the Workshop~3 agents, the VAE caught \monVaeCaught{} of the failures
without labels, but it also flagged harmless novelty. When labels exist, the supervised models of Workshops~1 and~2 should be
preferred.

\section*{Reproducibility}
Code in \path{src/w4/}: \path{autoencoder.py} (models, losses, training loop),
\path{train_ae.py} (experiment plan and logging), \path{evaluate_ae.py}
(reconstruction, anomaly detection, denoising, generation), \path{analyze_latent.py}
(latent analysis) and \path{make_report_w4.py} (tables and numbers of this report).
The frozen Workshop~2 model is loaded by \path{src/w4/w2_encoder.py} with the
Workshop~2 code (\path{src/w2/model.py}, \path{src/w2/sequences.py}). The
Workshop~3 monitoring is run by \path{src/w4/monitor_w3.py} with the Workshop~3 code
in \path{src/w3/} and the agent \path{checkpoints/workshop3/agent_best.pt}.
Checkpoints are in \path{checkpoints/workshop4/} (\path{ae_best.pt},
\path{vae_best.pt}, the normal-rows models, the supervised control and
\path{w2_lstm_best.pt}) and logs in \path{runs/workshop4/}. Seed 42; split
\path{data/splits.json}; scikit-learn 1.6.1.

\begin{thebibliography}{99}
\bibitem{barto1983} A.~G. Barto, R.~S. Sutton, and C.~W. Anderson, ``Neuronlike adaptive elements that can solve difficult learning control problems,'' \emph{IEEE Trans. Syst., Man, Cybern.}, vol.~13, no.~5, pp.~834--846, 1983.
\bibitem{brockman2016} G.~Brockman \emph{et al.}, ``OpenAI Gym,'' arXiv:1606.01540, 2016.
\bibitem{hinton2006} G.~E. Hinton and R.~R. Salakhutdinov, ``Reducing the dimensionality of data with neural networks,'' \emph{Science}, vol.~313, no.~5786, pp.~504--507, 2006.
\bibitem{vincent2010} P.~Vincent, H.~Larochelle, I.~Lajoie, Y.~Bengio, and P.-A. Manzagol, ``Stacked denoising autoencoders: Learning useful representations in a deep network with a local denoising criterion,'' \emph{J. Mach. Learn. Res.}, vol.~11, pp.~3371--3408, 2010.
\bibitem{goodfellow2016} I.~Goodfellow, Y.~Bengio, and A.~Courville, \emph{Deep Learning}. MIT Press, 2016, ch.~14.
\bibitem{kingma2014} D.~P. Kingma and M.~Welling, ``Auto-encoding variational Bayes,'' in \emph{Proc. ICLR}, 2014.
\bibitem{doersch2016} C.~Doersch, ``Tutorial on variational autoencoders,'' arXiv:1606.05908, 2016.
\bibitem{higgins2017} I.~Higgins \emph{et al.}, ``$\beta$-VAE: Learning basic visual concepts with a constrained variational framework,'' in \emph{Proc. ICLR}, 2017.
\bibitem{bowman2016} S.~R. Bowman \emph{et al.}, ``Generating sentences from a continuous space,'' in \emph{Proc. CoNLL}, 2016, pp.~10--21.
\bibitem{burda2016} Y.~Burda, R.~Grosse, and R.~Salakhutdinov, ``Importance weighted autoencoders,'' in \emph{Proc. ICLR}, 2016.
\bibitem{chalapathy2019} R.~Chalapathy and S.~Chawla, ``Deep learning for anomaly detection: A survey,'' arXiv:1901.03407, 2019.
\bibitem{schlegl2017} T.~Schlegl, P.~Seeb\"ock, S.~M. Waldstein, U.~Schmidt-Erfurth, and G.~Langs, ``Unsupervised anomaly detection with generative adversarial networks to guide marker discovery,'' in \emph{Proc. IPMI}, 2017, pp.~146--157.
\bibitem{kingma2015} D.~P. Kingma and J.~Ba, ``Adam: A method for stochastic optimization,'' in \emph{Proc. ICLR}, 2015.
\bibitem{rousseeuw1987} P.~J. Rousseeuw, ``Silhouettes: A graphical aid to the interpretation and validation of cluster analysis,'' \emph{J. Comput. Appl. Math.}, vol.~20, pp.~53--65, 1987.
\bibitem{alain2016} G.~Alain and Y.~Bengio, ``Understanding intermediate layers using linear classifier probes,'' arXiv:1610.01644, 2016.
\bibitem{vandermaaten2008} L.~van der Maaten and G.~Hinton, ``Visualizing data using t-SNE,'' \emph{J. Mach. Learn. Res.}, vol.~9, pp.~2579--2605, 2008.
\bibitem{mcinnes2018} L.~McInnes, J.~Healy, and J.~Melville, ``UMAP: Uniform manifold approximation and projection for dimension reduction,'' arXiv:1802.03426, 2018.
\end{thebibliography}

\end{document}
